% =====================================================================
%  R. Nielsen, ON FATE AND PROVIDENCE.  Copenhagen: Otto Schwartz, 1853.  83 pp.
%  (Om Skjæbne og Forsyn; reissued unchanged in 1867 as »Om Tilværelsens ideale Magter: Skjæbne og Forsyn«.)
%  ENGLISH TRANSLATION, from transcription.tex (Danish).
%  GENERATED by ebook-method/skjaebne-og-forsyn/tr_build.py --write from this template and the two translators' parts in
%  texts/nielsen/skjaebne-og-forsyn/.parts/tr/ (TRANSLATOR-BRIEF.md; terms: GLOSSARY.md, approved 2026-10-07).
%  Conventions: TRANSLATION-PLAYBOOK.md.  Page numbers of the 1853 print in the margin.  Emphasis as in the Danish:
%  \textbf = the print's bold Fraktur, \emph = its letter-spacing.  Scripture: Nielsen's Danish, in King James wording
%  where his Danish matches it.
% =====================================================================
\documentclass[12pt,a4paper,openany]{book}
\usepackage[utf8]{inputenc}
\usepackage[T1]{fontenc}
\usepackage{libertinus}
\usepackage{libertinust1math}
\usepackage{textalpha}
\usepackage{graphicx}
\usepackage[english]{babel}
\usepackage{enumitem}
\usepackage[protrusion=true,expansion=true]{microtype}
\usepackage{setspace}
\usepackage[margin=1.3in]{geometry}
\usepackage{fancyhdr}
\usepackage{hyperref}
\hypersetup{pdftitle={On Fate and Providence}, pdfauthor={Rasmus Nielsen}, hidelinks}
\setlength{\headheight}{15pt}
\pagestyle{fancy}
\fancyhf{}
\fancyhead[LE,RO]{\thepage}
\fancyhead[C]{\textit{R. Nielsen: On Fate and Providence (1853)}}
\renewcommand{\thefootnote}{\ifcase\value{footnote}\or *)\or **)\or ***)\else
  \arabic{footnote})\fi}
\setstretch{1.3}
\usepackage{marginnote}
\newcommand{\opage}[1]{\marginnote{\footnotesize\textit{[#1]}}}
% headings as in the Danish edition (the same macros as transcription.tex)
\newcommand{\overskrift}[1]{\begin{center}{\Large #1}\\[0.3em]\rule{3em}{0.4pt}\end{center}\addcontentsline{toc}{section}{#1}}
\newcommand{\afdeling}[3][]{\begin{center}\ifx&#1&\else#1\\[0.6em]\fi\textbf{#2}\\[0.3em]{\Large #3}\\[0.3em]\rule{3em}{0.4pt}\end{center}\addcontentsline{toc}{section}{#2 #3}}
\newcommand{\delafsnit}[1]{\begin{center}{\Large\textbf{#1}}\end{center}\addcontentsline{toc}{subsection}{#1}}
\newcommand{\stykke}[2]{\begin{center}\ifx&#1&\textbf{#2}\else\textbf{#1}\ifx&#2&\else\\[0.2em]\textbf{#2}\fi\fi\end{center}}
\newcommand{\motto}[1]{\begin{quote}\textbf{#1}\end{quote}}

\begin{document}

% PDF 3, the title page, p. [1]
\begin{titlepage}
\begin{center}
\vspace*{4em}
{\Large On}\\[1em]
{\Huge Fate and Providence}\\[1.5em]
by\\[1em]
{\Large\textbf{R. Nielsen,}}\\[0.5em]
{\small Professor of Philosophy at the University of Copenhagen.}\\[3em]
\rule{8em}{0.4pt}\\[3em]
\textbf{Copenhagen.}\\[0.5em]
Published by Otto Schwartz.\\[0.5em]
{\small Berling's Printing House.}\\[0.5em]
\textit{1853.}\\[4em]
{\small Translated from the Danish}
\end{center}
\end{titlepage}

% ===== TEXT (pp. 3--83) =====
% --- p. 3 ---
\setcounter{footnote}{0}%
\opage{3}\overskrift{Introduction.}

It is impossible, by an ordinary consideration of historical
events, to recognize the difference between fate and providence.

When Jerusalem was destroyed, the Jew said:
``it is Jehovah, who avenges the transgression of his law'';
the Nazarene said: ``it is the blood of the Crucified One
that comes upon the head of the people''; but the Roman said:
``it is fate.''

As the Jew, the Nazarene, and the Roman then
judged, each according to his mood and his standpoint,
so to this very day every party and
every human being judges, according to mood and standpoint,
of every historical event. It does not, then,
depend on the event, but on the mood and
the standpoint.

It is impossible, by an ordinary consideration of
the particular remarkable occurrences in the lives of particular human beings, to make a distinction between fate and providence.

Two human beings can outwardly experience the
same incidents, come into the same dangers, be saved
% --- p. 4 ---
\setcounter{footnote}{0}%
\opage{4}in the same unexpected and wonderful way: the one
sees in all this a play of fate, the other sees
a providence. It does not, then, depend at all on
what it is that is experienced, but on \emph{how} it is experienced;
it does not depend on the occurrence, but on
the way of thinking.

It is impossible for a human being to have his faith
in providence directly confirmed by the events
that meet him; for he must, after all, interpret all events
according to his faith. If prosperity comes: it is providence;
if adversity comes: it is also providence; if everything in
existence is comprehensible and happy, and the future bright:
it is providence; if everything in existence is confused and
troublesome, and the future dark: it is also providence.

He who does not thus in everything subordinate his own
human judgment to his faith in providence does not
believe in providence at all.

But if events, in history as well as
in the lives of particular human beings, admit of being interpreted in such
widely different ways, then, as regards fate
and providence, we can learn nothing, absolutely nothing, from the
course of events.

% --- p. 5 ---
\setcounter{footnote}{0}%
\opage{5}\afdeling{A.}{On Fate.}
\stykke{I.}{What Is Fate?}

The cold, inflexible, and yet so capricious,
arbitrary power which we call fate is essentially
grounded in the interconnection of natural things,
in the law-governed order of existence,
that is, in the world-reason itself.

The law of reason and the necessity conditioned by it
are, however, in and for themselves by no
means yet to be called fate. Insofar as the
things that exist immediately subordinate themselves to the
universally valid laws, necessity, the
outer as well as the inner, is to be regarded as the ground of reciprocal action,
of wholeness, and of harmony.
The course of the stars, the change of the seasons,
% --- p. 6 ---
\setcounter{footnote}{0}%
\opage{6}the growth of plants, all such things are grounded in the interconnection of nature,
and insofar in necessity;
but here there is no fate.

Only in the animal world, where individual
life expresses itself with self-awareness and selfish
purpose, does fate begin to dawn. What
an animal cannot suffer; how various, in
this respect, is the fate of animals: the caged
bird and the free bird each have their own fate;
and then the fate of domestic animals, --- if horses
and asses could speak!

But horses and asses cannot speak; this
imperfection mitigates their fate. The animal's
dependence on instinct, its lostness
in the moment, the dimness of its consciousness, which
veils the possibilities of the future and hides the
past in oblivion, prevents it from
setting itself up as a free end in itself, bends it gently
under the yoke of natural necessity, and is like
a covering that deadens the blows of fate.

Only where there is a self-purpose, a plan
of the will, a personal goal, can there be
talk of a power that cows the purpose, thwarts
the plan, and disrupts the goal. Among all earthly creatures, therefore, man is the only one  % print: Jorgens (misprint)
% --- p. 7 ---
\setcounter{footnote}{0}%
\opage{7}to whom nature has granted the great, ambiguous
privilege that, in the proper sense, so long as
he stands on the ground of worldliness,
he is to stand under fate. Rich endowments, mighty
will, bold spirit, all the natural gifts of
personality cannot save him from the
dark, enigmatic power; for the greater the personality,
the mightier fate. Fate is the enemy of all that is personal;
therefore the gods must feel its power, therefore
Jupiter with his thunderbolt must bow before Fatum.

\stykke{II.}{The Dominion of Fate.}

If it is a distinguishing mark of fate that it
everywhere confronts personality with hostility,
then it follows that its dominion
must be as extensive as personal
striving.

In the acquisition of temporal goods, in the sphere
of ideas, indeed even in the moral, we must
distinguish between that which depends on us and
that which does not depend on us. If now that
which does not depend on us took account of our
personal need, and came to the aid of our rational will,
% --- p. 8 ---
\setcounter{footnote}{0}%
\opage{8}then it would not be fate, but a
higher wisdom governing all. Viewed in worldly terms,
the reverse is just the case. The universal
superior power on which our striving depends is
so far from heeding the personal in us that
it rather behaves indifferently, indeed
even mockingly toward it: this is precisely the fatal thing.

If we consider the difference between the healthy and the
sick, the rich and the poor, the high and the low in the world,
it is seen at first glance that the distribution
of temporal goods by no means takes place
according to need and merit, but that the distributing world-power has followed arbitrariness,
has let fortune and misfortune, good luck and bad luck rule
the division, and, in the guise of good luck as well as
bad luck, has made a fool of human personality.

If we turn from the outer to the inner, a similar arbitrariness
repeats itself.
It is a human privilege to be able to
rise above the particular, the petty,
to set ourselves universally valid goals, that is to say:
to live for ideas. But this privilege is not
equal for all; here we meet the fateful
difference between the gifted and the
% --- p. 9 ---
\setcounter{footnote}{0}%
\opage{9}ungifted, the spirited and the spiritless. The
expression, to live for an idea, is moreover ambiguous
so long as it is left undetermined to which
sphere the idea belongs. There are personal
and impersonal ideas.

The ideas of art, science, and the state are
impersonal; the idea of morality is personal.

To live for an impersonal idea is to give
oneself up to the impersonal. Every spiritual devotion
to the spiritual heightens the self-feeling; therefore
geniuses and the inspired have such a powerful
intense life in themselves. But the impersonal
finally gains the superior power. The bond of
inwardness that binds the distinguished to the idea and
makes them its powerful instruments is at the same time
the chain in which fate catches personality  % print: Personligligheden (misprint)
and binds it to its wheel: the work
is completed, the master perishes; the idea
triumphs, its instrument lies shattered (a
Columbus, a Kepler, a Mozart, a Fulton
could, each in his own language, give an explanation).
Hard is fate when it arbitrarily distributes
the finite goods, plays with the fortunate
and torments the unfortunate; but cruel is it
when, as impersonal idea, it steals into
% --- p. 10 ---
\setcounter{footnote}{0}%
\opage{10}the inner being of man and sucks the blood of personality.

In order to withstand fate, man ought therefore
to submit himself to the moral world-order
and practice virtue, for the idea of virtue is
personal. This counsel is old and good,
but, understood in worldly terms, insufficient; as far
as the laws of nature extend, so far extends also fate's
realm.

As we distinguish, with regard to the impersonal
ideas, between the gifted and the ungifted,
so we must, with regard to the moral
idea, make a distinction between noble and ignoble natures.
The moral natural dispositions seem, namely,
to be as arbitrarily distributed as the talents
and the temporal goods. There are individuals
who cannot but be pious,
virtuous, lovable, because it lies in their nature;
and there are individuals who by nature
are as good as predestined to crime.
Certainly man has what is called a free
will; but what can the will do against passion?
what becomes in reality of
the moral will, where the natural ground
is lacking? Half-heartedness, hypocrisy, quackery,
% --- p. 11 ---
\setcounter{footnote}{0}%
\opage{11}mediocrity! The will grips powerfully only when it is
itself powerfully gripped, but it is gripped powerfully only
by that which dawns on the inner sense,
and satisfies personally. If now the moral
idea has its validity solely in personal
self-satisfaction, then he who
satisfies himself in what is bad, when he really
satisfies himself, is just as justified as he
who satisfies himself in what is good; the difference
is set by fate and is erased again by fate;
for death makes us all equal: death is
the last fate.

\stykke{III.}{The Struggle against Fate.}

The worldly universal spirit is an impersonal
spirit that would devour the human self;
worldly personality is at bottom an
egoistic personality, which in defiance of
the universal would assert the human self. Between
the impersonal universal and the egoistic
personal the relation is strained. The superior power
of the impersonal in the relation is fate; therefore
personality struggles against it, and maintains itself,
as far as it is able, in the struggle against fate.

% --- p. 12 ---
\setcounter{footnote}{0}%
\opage{12}According to the particular personal demand
and the way in which the demand is carried through,
we characterize the struggle as: the \emph{civil, the heroic, and the moral}.

The civil world is the world of finitude,
of temporal goods. In a world built on
sheer conditions, the rational
will cannot then think of making
any unconditional demand or requiring unconditional
satisfaction. The finite goal is a tolerable
condition: personality bows to circumstances, and comes to terms with fate.
But the agreement that is thus concluded
between the egoistic and the universal is nevertheless
transient, unreliable, and false; the contracting
parties are essentially each other's enemies:
the tolerable condition soon becomes intolerable, and
personality is driven to despair. The
mighty oppress the lowly, and the lowly hate
the mighty; the rich despise the poor,
and the poor envy the rich; good fortune corrupts the fortunate, misfortune the unfortunate; corruption
begets bitterness, bitterness breaks
out in rebellion.

When Cadmus sowed the dragon's teeth, there
% --- p. 13 ---
\setcounter{footnote}{0}%
\opage{13}sprang up armored men, who slew one another:
these armored men are the sons of fate.

Seen in the phenomenon, it is the rebellion of
human beings against one another; seen
essentially, it is the rebellion of finite, egoistic personality
against fate. Seen in the phenomenon,
it is an altered condition, a new
order of things, that emerges from the struggle;
seen essentially, it is the same condition, the
same transient, unreliable, false agreement
between the personal and the impersonal, the egoistic and the universal, thus
still as before: the same fate.

In the civil tension between the selfish
and the universal, personality must at once
feel itself the weaker; which is why it also
moves between the two extremes: pliant
submissiveness and desperate insolence. If
personality, on the other hand, feels itself as strong as
the universal, the tension becomes passionately
spiritual, and the struggle \emph{heroic}.

The heroic will is not content to receive
conditions, it would itself prescribe conditions;
it does not make do with a tolerable
% --- p. 14 ---
\setcounter{footnote}{0}%
\opage{14}condition, it demands unconditional satisfaction.
Here there is moreover a higher intermediary, in which
the strained extremes, personality and fate,
unite; this intermediary is the idea.
The heroic will wills an idea, and insofar
the universal; it also wills this idea realized
in a determinate worldly goal, thus in a
finitude: here fate lurks. By marking out
its goal for itself, personality measures itself
against fate. That fate has a superior power
is at once evident; but instead of giving in,
the will nevertheless prescribes to this superior power
its conditions, in that it prescribes to itself a
choice. This choice is between All and Nothing,
the attainment of the goal and the downfall of personality.

``Whatever it cost, I go to the
goal!'' with this watchword freedom has raised
itself above fate; so a Croesus interprets the oracle,
so a Wallenstein reads in the stars, so
a Caesar crosses the Rubicon, and so begins
the struggle.

It is the attainment of the goal that gives the
personal its worth; therein lies ``that, when
the goal is taken away, personality is in itself
% --- p. 15 ---
\setcounter{footnote}{0}%
\opage{15}something wholly indifferent; the heroic resolution
thus expresses the contradiction that the
personal proves its highest validity by
setting itself as a means for the goal, as that which is in
itself indifferent. If this contradiction
dawned on consciousness in a clear and sudden
thought, it would be like a dagger thrust
that instantly killed the hero; but the contradiction
is veiled by the idea, is sensed as unrest in the mind, and transforms itself into a driving wheel
of events. Permeated by the idea, the
personal self-feeling is lost in the impersonal.
The hero is most himself when he is
most outside himself in his irresistible
longing toward the goal; seized by an inexplicable
power, he feels himself powerful; this feeling of power
is his initiation, his mystical covenant with
fate.

It is the cunning of fate to catch the personal
by personality itself, as it is
the pleasure of fate to trample its captive underfoot.
Fate beguiles in the form of good fortune; it paralyzes
as the riddle that awakens dread; it strikes
in the shape of Nemesis; the beguiled Croesus,
the paralyzed Wallenstein, the crushed Caesar.

% --- p. 16 ---
\setcounter{footnote}{0}%
\opage{16}In the heroic struggle the fusion of the personal
with the impersonal is a mystery.
To unveil this mystery is a task
for ``the wise man.'' Personality turns
inward upon itself, in order to clarify and realize
its free determination; thus the
struggle becomes moral.

In order to measure itself against fate, personality must
first of all give an account to itself of what
it is that really stands in our power, and
what all that signifies which does not stand in our power.

If it is now seen that the distribution of temporal
goods and evils: health and sickness,
riches and poverty, civil freedom and civil
bondage, etc., does not essentially stand in
our power, then liberation must begin
with this, that the difference between temporal goods
and temporal evils is reduced to a difference indifferent
in itself.

To ``the wise man,'' then, it is wholly indifferent
whether he is healthy or sick, rich or poor,
master or slave. Certainly it is more pleasant
to be healthy than to be sick; but since the will
can prove its strength just as well in sickness
as in health, the difference between the
% --- p. 17 ---
\setcounter{footnote}{0}%
\opage{17}pleasant and the unpleasant cannot touch the personal
self.

To put up with the unpleasant and to accommodate oneself
to adversity is resignation.

Already the civil struggle makes resignation
necessary. He who craves temporal
goods must put up with temporal losses; but it
is only finite resignation where the loss of the
one is compensated by the attainment of the other.
Thus in finitude fear and hope, sorrows
and joys, loss and compensation alternate; even the loss
of loved ones, however painful it may be at the moment,
can be overcome, compensated --- with time;
for time heals all wounds.

This consolation of finitude, this temporal
alternation of fear and hope, of the misery of the moment
and the jubilation of the moment, ``the wise man'' despises.
Raised in calm, proud self-feeling above the changeable,
he resigns, not now and then with
regard to this or that finite thing, but once
for all with regard to everything finite, in that he recognizes
that the temporal by its concept is and
must be temporal, the perishable perishable.

``If you have'' --- says the wise man --- ``a beautiful
% --- p. 18 ---
\setcounter{footnote}{0}%
\opage{18}earthen vessel, remember well that it can be broken.
If you forget this, it will painfully
disturb you when it happens; but if you
always remember it, chance cannot surprise
you. If you love your wife and your son, then
bear in mind that they are mortal. If you forget this,
the pain may one day confuse your mind; but
if you always remember it, you will calmly, without
sigh or lament, follow them to the grave.''

What, then, is fate? In its phenomena
an arbitrary play of prosperity and adversity,
good luck and bad luck, but according to its essence that
which is grounded in the nature of things, the necessary,
the universally rational. From this is derived the moral
law, that the ethical personality must
will nothing other than what reason wills, and
wish for nothing other than what follows from the course of
things.

``You shall not weep when you sit in
the theater and see a tragedy; for it is not
sorrowful that the hero falls, if it is necessary.  % print: ørgeligt (misprint)
What must happen ought to happen, and
with that you shall be content.''

Devoting himself to the service of necessity,
% --- p. 19 ---
\setcounter{footnote}{0}%
\opage{19}will the personal self transform necessity
into freedom and make freedom irresistible
as necessity. Thus reconciled with
the laws of existence, the wise man says: ``lead me,
O Fate!''

In order to act in the world the will must set itself
a worldly goal; but here is the difference between
``the wise man'' and ``the fool.'' The fool makes
his personality a means for the goal; the
wise man makes every finite goal a means for
his personality.

``The world,'' says the wise man, ``is a great theater.
Human beings are actors. The Deity
is the poet who has the play performed, and
distributes the roles. So the one gets the king's,
the other the beggar's, the one the master's,
the other the slave's role. It is not for the actor
to reject his role, but
to play his role. Only he who plays his
role well wins the applause of the Deity.''

The Deity whose applause the wise man craves
is, in the proper sense, personified
Fate; the moral self is the personification
of Fate. Infinite resignation is
% --- p. 20 ---
\setcounter{footnote}{0}%
\opage{20}thus after all an infinite egoism; it is the actor
who on his own authority has composed for himself
a role; it is the actor who plays for
himself, and pleases himself in the role; it is the
actor who in the end reaps his
own applause.

That this is so without exaggeration
is plainly seen in the last act.

The fateful personality, which here
has made itself its own deity, cannot
prove its self-will more plainly than by this,
that it is master of its own life, and therefore itself
chooses and determines its own death. The personal
life is finite; the finite is indifferent.
With freedom and deliberation the friend of Fate
carries out what Fate would otherwise
force upon him: he takes his own life. Personality
asserts itself by annihilating
itself. Herewith the moral struggle and
the moral drama end.

% --- p. 21 ---
\setcounter{footnote}{0}%
\opage{21}\afdeling{B.}{On Providence.}
\stykke{}{What Is Man?}

Yes, what is man? Something great and
glorious; something paltry and wretched! There is
on earth no creature that can become as
happy as a human being; for there is on earth no creature that can become as unhappy
as a human being; none that can rejoice as
a human being, because none can grieve as a
human being; none that can fear what a human being
can fear, because there is none that can
hope as a human being can hope; none that
can hate as a human being, because there is none
that can love as a human being; none that can
be arrogant as a human being; none that can
humble himself as a human being; none that can
mock as a human being; none that can worship
as a human being; none, none that can suffer
the wrong a human being can suffer; none, none
that can do the wrong a human being can do
to himself and others. What, then, is man?

% --- p. 22 ---
\setcounter{footnote}{0}%
\opage{22}All the helpless who are in the power of nature
are helped by nature; man is in the power of
fate: who helps man?

\delafsnit{Providence Revealed.}
\stykke{I.}{}
\motto{For he maketh his sun to rise on the evil and
on the good, and sendeth rain on the just and on the
unjust.}

But is that then a help? Can the
good man console himself that the sun rises also on
the evil? Can the just man rejoice that
heaven blesses the unjust? We know that
it goes thus in the world; we know that the laws of
nature recognize no difference between evil and
good, between just and unjust. The storm
rises, and the ship sinks, and the human beings
perish; they perish, both evil and good,
both just and unjust. The earth opens
and swallows the city, and the human beings
are destroyed; they are destroyed, both evil and good,
both just and unjust. The plague passes
through the lands, and human beings die of
the deadly sickness; they die, both evil and
% --- p. 23 ---
\setcounter{footnote}{0}%
\opage{23}good, both just and unjust. Can
such a thing then signify a providence? It looks very much like
fate.

Yes, if everything personal is finite, and
the eternal itself impersonal, then it is fate.
But if the personal is at the same time the eternal,
if true personality is the truth
of the eternal: then there is a providence.

If the eternal personality were not itself
human, there would be for man no providence
and no help; if man of himself
were the eternal personality, then every
human being would be his own providence and in the eternal sense
his own helper.

Every human being who must confess: ``I am
not the truth, but I need the truth
personally,'' must relate himself in faith to providence,
and cannot be his own providence.

Only the human being who dared to say of himself:
``Before Abraham was, I am, the truth
of the eternal, the true personality,'' only such a
human being can relate himself knowingly to providence,
can as a human individual be under
providence, and yet be Providence itself, who makes
himself the object of faith.

% --- p. 24 ---
\setcounter{footnote}{0}%
\opage{24}The Son of Man, who knew within
himself that he was the true providence revealed in flesh and blood,
and therefore made himself the object of faith,
has accordingly not taught us
to see the justice of providence in a finite
distribution of temporal goods and temporal evils.

To pay for the temporal with the eternal, the
eternal with the temporal, is for all eternity unjust;
but to pay for the temporal with the
temporal, the eternal with the eternal, that is like
for like, that is divine justice.

If it were really the thought of Providence
that good and evil should be requited
temporally, should we not then expect that he who
became man solely in order to be
the revealed Providence would receive riches
and power and all other temporal goods? Since
this now does not happen: what can be the reason?
Might Providence perhaps envy a human being
the temporal? Or might he who is the
true man set a kind of honor in doing without
and despising what is human? That would be
a mockery of everything divine and human.
When he, then, who could have all the glory of the
world, does not have anything on which he can
% --- p. 25 ---
\setcounter{footnote}{0}%
\opage{25}lay his head, this abasement is not
to be ascribed to an arbitrariness of the will, any more
than to a necessity of fate. It is human
self-denial for the sake of truth,
it is the infinite self-affirmation of personality,
which in this lowliness reveals the
divine superiority, the perfect justice.

For millennia now humanity has struggled
for temporal requital; and still the sun rises on the evil and the good, and still it rains
on the just and the unjust. He who
looks at the course of the world in worldly fashion sees in all
this only fate; but he who looks in faith
at the Son of Man sees the will of the Father,
who is in heaven.

If it were the thought of Providence that everyone
who struggles for a universally valid idea
should be recognized by human beings, and be vindicated
in proportion to the worth of the idea and the power of
sacrifice with which he struggles: should we not then
expect that He who not only struggled for the highest
idea, but was the idea itself, would have won
for himself the highest recognition in the world?

All the great misunderstood, who otherwise have
% --- p. 26 ---
\setcounter{footnote}{0}%
\opage{26}struggled for ideas, and have been by their contemporaries
rewarded with ingratitude, hatred, and contempt, could yet
demand only a conditional recognition; it was
yet only relative ideas for which they struggled; and
even if one struggled for the highest, he could
yet not say with truth: ``I, who struggle
for the highest, am myself the highest.''

The Son of Man, who struggles for the eternal
truth, bears witness of himself: ``I am the truth'';
the Son of Man can therefore make himself
the object of faith, can demand unconditional
submission, divine reverence.

All the great misunderstood, who otherwise can
complain of human ingratitude, must
yet admit that their contemporaries, by denying
them recognition in the name of the idea, harm
only themselves to a certain degree; the idea after all
comes sometime to the benefit of the race: the individual's
personality is indifferent.

The Son of Man, who is personality
itself, and must demand unconditional recognition for
the sake of personality, knows with irrefutable
and painful certainty that every human being who
denies him what he demands inflicts upon himself
an eternal personal wrong.

% --- p. 27 ---
\setcounter{footnote}{0}%
\opage{27}All the great misunderstood, who otherwise were put in the wrong
although they were in the right, could with reason
complain of fate; and we, who now see the
wrong they had to suffer, we can lament them, indeed
we can weep over them, when we think of
their fate. But He, the One, who was put in the
wrong, not although but because he was in the
right, eternally, unconditionally in the right, he will not be lamented,
not pitied, not wept over; he says expressly:
``Weep not over me, but weep
over your own and your children's fate.''

So the just one was among the unjust;
the sun rose on the evil and the good, and
as the sun sank, darkness rose on the good and
the evil; and the darkness strengthened the evil, and in an
hour of darkness the just one suffered for the
unjust. It looks like a fate;
yet it is not fate, it is the Father's will,
who is in heaven.

If it were the thought of Providence that human beings
themselves should find the way by the light of
reason, then the Son of Man would never have made
himself the object of faith.

Only like is recognized by like, the
personal by the personal, but the personalities of God and
% --- p. 28 ---
\setcounter{footnote}{0}%
\opage{28}man are not each other's
equals.

``Honor your friend, despise your enemy; procure for yourself
justice in time, if you can, for in eternity no
legal proceedings are held; good and evil, right
and wrong are only differences of finitude, personal
oppositions, which are all dissolved in the
eternal, impersonal One.'' Here is the wisdom
that denies Providence and deifies fate.

``Love your friend, hate your enemy; for your
friends are the friends of the Eternal, your enemies
the enemies of the Eternal: see, the sword is in your
hand, that the justice of heaven may be fulfilled
on earth.'' Here is the passion
that would transform Providence itself into a
fate.

The imagined superior power of personality is in
both forms an eternal impotence. He who despises
his enemy essentially despises himself;
for he despises the personal: the superior power
of contempt is an impotence of egoism. He
who makes his enemy God's enemy rages
against God: the fanatical superior power is a frenzied
impotence.

In order to save human personality
% --- p. 29 ---
\setcounter{footnote}{0}%
\opage{29}from its natural impotence, Providence itself came in
the form of a servant. Had he come in the form
of a ruler with outward divine
superior power, we should have remained in our natural
impotence. But he denied himself, let the
world keep the outward superior power, and
became obedient unto death, that in self-denial
and obedience he might reveal the infinite
superiority of personality. His enemies
were certainly God's enemies; he saw their
blindness, he grieved over the hardness of their
hearts; but he did not despise them, he  % print: men. (misprint)
did not curse them: he prayed for them, and
died.

Ye have heard that it hath been said: ``Thou shalt
love thy friend, and hate thine enemy'';
but I say unto you: ``Love your enemies,
bless them that curse you,
do good to them that hate you,
and pray for them which despitefully use and persecute you, that ye may be
the children of your Father who is in the heavens.
\emph{For he maketh his sun to rise on the
evil and on the good, and sendeth rain on the
just and on the unjust}.

% --- p. 30 ---
\setcounter{footnote}{0}%
\opage{30}\stykke{II.}{}
\motto{Yea, even all the hairs of your head are numbered; fear ye
not therefore.}

But are these words not to be understood as
a poetic exaggeration? If divine
justice is really exalted above
every temporal difference between good and evil,
is it then conceivable that the Eternal should sink
so deep into the temporal and be so exacting
about the particular? Is it conceivable that
the Providence of mankind should count all the
children of men, as a shepherd counts his sheep?
Is it worthy of the All-Knowing to count so pettily,
even the sparrows, even the hairs on the
head of man?

``All the hairs of your head are numbered, fear ye
therefore \textbf{not}!'' If these words were only human
words, we should have to regard them
as an exaggeration in an outburst of enthusiasm;
for in the world of dangers that lies
between man and God, a human being
has much to fear. But when He who
speaks the Word is the Word itself, then we must learn
from him where the danger is, and what we have to
fear.

% --- p. 31 ---
\setcounter{footnote}{0}%
\opage{31}There is in the world of dangers only one unconditional
danger, only One that is fearful, only One
who is to be feared. Therefore the Word, the
revealed Providence, says: ``I will show you whom
ye shall fear.''

And Providence is human. He who
speaks of the danger, and shows us whom we shall
fear; He who surveys all the dangers of
finitude and says: ``fear not!'' he knows what
danger is; his life was familiar with danger.
Danger sought him in his cradle; it followed
him to his grave. Everywhere, in the desert as
in the cities, on the plains of Galilee, among
the mountains of Judea, everywhere he went, danger
went after him, and lay in wait for him.

Watchful as Providence, courageous as he
who is under Providence, the Son of Man goes
along the way of dangers, and fears nothing in the world.
He works, and nothing stops him;
he goes to his goal, and no one lays a hand
on him until his hour comes.

But when his hour had come he went
voluntarily to meet danger; it did not surprise him
as an accident, it did not come upon him as
a fate; it came because he let it come.
% --- p. 32 ---
\setcounter{footnote}{0}%
\opage{32}Nor did it come in outward uproar, not in
a storm of human passions, not in a
clamor of hostile weapons; it came in all
stillness, alone to the Alone: only thus
did he let it come.

Alone that night in the quiet garden, he let
danger come to him, and struggled with it, not
as the strong struggle, but as the frail, so that in flesh and blood there should not be
a difference of fate between the strong and the
frail, but a victory of devotion in
frailty.

Alone that night in the quiet garden, he let
danger pass as a mortal anguish through his
soul; he did not defy it, he did not despise it,
he did not proudly lift his head like him
who sacrifices himself to fate; he fell on his
face, and bowed himself humbly under the Father's
will.

Alone that night in the quiet garden, he let all
the dangers of finitude transform themselves in inwardness
into one danger. So he struggled in inwardness,
and conquered in inwardness; and when he then rose
from the struggle, behold, then it was over with
danger, then there was no danger at all.

% --- p. 33 ---
\setcounter{footnote}{0}%
\opage{33}Yet --- what do I say? Is it really over
with danger: now the real danger first comes,
for now the guard comes with swords and
staves. Yes, now reality comes, but
--- the danger has turned. He who asks:
``Whom seek ye?'' is undismayed; the guard
falls to the ground in dismay.

The night when the Son of Man was betrayed,
scourged, and condemned, is in reality
a night of dangers; but the danger has turned.
Here is danger for a Judas; here is danger for
a Caiaphas. Here is danger for a Pilate,
who washes his hands; here is danger for
the people, that the blood of the Innocent comes upon
their head. Yes, here are many dangers and for
many; but for Him, the One, for him there is
no danger.

Were there danger for him, he would surely get relief;
were his life, his cause really in
danger, his servants would surely fight for
him, and more than twelve legions of angels arm
themselves for his defense. Now he stands without defenders, weaponless, alone among his enemies:
why? Because the danger is already over; because his
hairs are numbered.

% --- p. 34 ---
\setcounter{footnote}{0}%
\opage{34}``Not a sparrow falls to the ground without the Father's
will; yea, even the hairs of your head are
numbered.''

The world of finitude is a world of
contingencies. If now this or that
accidental thing is held fast in its accidental singularity, then
every such accidental singularity, as e.g.
that a sparrow falls to the ground, that a
hair falls from a human head, etc., is certainly
indifferent to God and man. Such a
holding fast of the accidental in its contingency
is, namely, spiritless, and the spiritless
is for all eternity indifferent.

The accidental first gains significance by
being seen in an essential relation to an Other
that is not accidental. This Other can
again be: either the rational ground of the world or
personality itself.

What contingency signifies in its connection
with the ground of the world and natural
causes is a scientific question, which
does not at all concern faith in providence.

What contingency, on the other hand, signifies as
a worldly matter between God and
man, is a personal question. This
% --- p. 35 ---
\setcounter{footnote}{0}%
\opage{35}question can be answered only by Him who  % print: Ham., (misprint)
reveals the thought of Providence.

Here is one born blind; let us hear the Master's
conversation with the disciples. ``And the disciples
asked him and said: Master, who hath sinned?
this man or his parents, that he is born
blind? And the Master said: neither this man
sinned, nor his parents; but that the
works of God should be made manifest
on him.''

The disciples and the Master are thus agreed in
the presupposition that chance cannot be
anything independent in relation to God, that Providence
in every accident removes the accidental.
But in that the disciples simply overlook chance,
they are entangled in a double misunderstanding, inasmuch as
they at once fail to recognize the reality of natural
conditions, and drag divine
justice down into the temporal. The Master,
who knows the thought of Providence, calls attention
to the difference between natural determination and
God's express will. The difference between
the sighted and those born blind is a natural difference in
the temporal, which does not directly correspond to God's
retributive justice. Insofar the accidental is thus
% --- p. 36 ---
\setcounter{footnote}{0}%
\opage{36}conceded. But the accidental
is not for God in its contingency;
it is determined as a means for his determinate
purpose. What was therefore the case with
that man born blind, that the works of God should
be made manifest on him, the same is just the case
with all those born blind and with all the possible
accidents in the world that concern the personal.

--- ``As long as I am in the world, I am the light
of the world.''

The natural man, who sees everything in the light of naturalness,
sees accident upon accident,
both in the course of the world and in his own life; it
is the congenital blindness in finite personality,
which cannot see Providence for sheer accidents.
The Son of Man, the personal truth,
opens the eyes of the blind, and thereby removes the
sad accident.

``As long as I am in the world, I am the light of the world.''

The history of Providence is a sacred history.
If sacred history were to be grasped directly,
like any other history, then
we would only need to follow the Son of Man on his
way, in order to see how he abolishes
% --- p. 37 ---
\setcounter{footnote}{0}%
\opage{37}contingency and reveals the work of God.
But in the sacred history of Providence the outer is
so dissimilar to the inner that only he who
sees everything in the light of faith sees Providence.

The richer the sources of help that personality has
in itself, the more easily it transforms the accident
into means and occasion; the more surprising
the accident is, the more surprising also is
the presence of mind that flashes through contingency.
Therefore genius is in league with
fortune, is itself fortune; therefore practical
genius counts as many victories as
decisive accidents.

If now the personal superiority
that reveals the thought of Providence were to be compared
with the superiority of genius: then we should
invite the spirited, the geniuses, to
go through the sacred history; then we should
say to these spirited, these geniuses: ``Look
for once rightly at the Son of Man; as he
masters contingency, no one on
earth has yet been able to master it; for to him nothing is
accidental, he abolishes and transforms all
accidents. Is it the manifold cases of illness;
he removes every accident, for he
% --- p. 38 ---
\setcounter{footnote}{0}%
\opage{38}heals all the sick. Is it the sad
cases of death; he removes them, when and
where he will, for he can raise the dead.
Is it one of malice's intriguing, ingeniously
contrived accidents, an accident like that with the tribute
money, or an accident like that with the woman
and her seven husbands in the resurrection; behold,
he masters all these contingencies. The lightning
of truth splits the web of malice, and the accident
falls like a fate on those who
harden themselves in untruth.

But although the personal omnipotence
in all these cases is unmistakable,
such a conception is in itself a blasphemy;
the Son of Man is not a hero of
fortune; he is the Holy One of Providence, the object of faith.
--- ``Believest thou that I can do this?'' ---

If you do not believe, the accident remains
in its sad contingency; but ``if
thou wouldest believe, thou shouldest see the glory
of God.''
--- ``As long as I am in the world, I am the light of the world.''

To the fate-born blind all finite
accidents are indifferent. With this everything
% --- p. 39 ---
\setcounter{footnote}{0}%
\opage{39}particular is dissolved in a uniform, impersonal, empty generality
that lacks the content of life and
the richness of colors. ``The wise man'' with all his
strength is therefore also like a spiritual shadow-image,
a specter of Fate, that dwindles into
the dim.

For him who ``believes in the Son of God,'' the
personal truth, the true personality, the
worldly accidental is a personal matter
between God and the individual human being.
Every accident that meets one is a particular task
which Providence sets for personal solution: as many
accidents, as many tasks. In the
finite solution possibility stands against
possibility, understanding against misunderstanding, danger
against danger: --- Man, do what you will,
but do it in unconditional faith! If you do
not believe, it will be for you a sad accident,
but --- ``if thou wouldest believe, thou
shouldest see the glory of God.''

--- ``As long as I am in the world, I am the light of the world.''

The outwardly changing manifold of
finite accidents is, when seen in the light of
truth, sheer particular conditions for the
% --- p. 40 ---
\setcounter{footnote}{0}%
\opage{40}development of the personal inner being. The accidents of life
are God's dispensations. Unbelief, which in
the dispensation sees only chance, falls to
fate; faith, which in every accident sees
the dispensation, is saved by Providence.

Man's relation to the Providence revealing itself in the sacred
history is therefore not
a uniform, colorless relation that loses itself
in impersonal generality; it is an infinitely
rich in content, manifold life-relation, nuanced to the smallest
details, between the personal
truth and human personalities.
Here in every scene faith stands against
unbelief, the sighted against the blind; here the
blind become seeing, the sighted blind; here
every event, every accident, every word,
every glance, indeed every moment has the significance
of infinity; here the light rises on evil and
good; here the dispensation comes upon the just
and the unjust; here the works of God are
made manifest; here the thoughts of many hearts are
made manifest.

--- ``As long as I am in the world, I am the light of the world.''

In the highest tension between
% --- p. 41 ---
\setcounter{footnote}{0}%
\opage{41}wisdom's plan, the arbitrariness of men and the world's
contingencies, Providence is revealed;
therefore the sacred history is a history of
sufferings.

In the history of suffering it looks as if
Fate itself would rise up to battle against Providence,
as if contingency were to rule, and human
arbitrariness tear itself loose from governance.

And in the history of suffering it looks as if
God's will had eternally predetermined each single event,
each single scene, each single trait.

--- So the One must ``die for the people,
that the whole people perish not.'' It
is a notion of malice, a despotic arbitrariness
of an individual in a single accident; and
yet, --- it is no arbitrariness, it
is an inspiration of the Spirit, it is a prophecy
that expresses the thought of Providence.
--- So the Just One is ``reckoned
among malefactors''; --- an outrageous, sad
arbitrariness, a sad accident! and yet,
it is no arbitrariness; it is the foreordained
counsel of wisdom, it is the fulfillment
of prophecy.

% --- p. 42 ---
\setcounter{footnote}{0}%
\opage{42}So the Crucified One dies: the sun is darkened,
the veil is rent, the graves burst\dots{}.
``marvelous occurrences, perhaps invented accidents''\dots{}.
Yes, but --- if thou wouldest
believe!

\emph{Are not two sparrows sold for a farthing,
and not one of them falls to the ground without your Father's will. Yea,
even all the hairs of your head are numbered.
Fear ye not therefore, ye are better than
many sparrows}.

\stykke{III.}{With God All Things Are Possible.}

But no miracles happen after all! ---
``Yes, that is just the disquieting thing, that a personality
who belongs to the age of wonders is here
presented as a model, that He who performs the
many wondrous deeds is just to be an example
for us, who live under the laws of nature. He who
can still the storm may well sleep in
the ship; but he who cannot still the storm
must watch and use his strength. He who
with a word can cast the guard to the ground can
easily put the sword into its sheath; but when
% --- p. 43 ---
\setcounter{footnote}{0}%
\opage{43}the word can do nothing, one must try the
naked steel.''

If these misgivings belonged merely to
declared unbelief, we should pay no heed
to them; but where is a faith so strong that it
does not struggle with doubt?

Most human beings surely know such
moments of sorrow, of anxiety, of tension,
in which the thought as if involuntarily forces itself
upon them: Oh, that a miracle might happen after all!

--- ``He hath saved others, himself he cannot
save.'' So sounded the mocking word
over Him who had helped so many, and
now was himself as one helpless.

All help in the temporal is in the thought of Providence
a means for the eternal: help can be given
by a miracle; but help can also be given
without a miracle. If the miracle is understood as
a help in the temporal, then it does not help at all;
it harms. If helplessness in the
temporal is understood as a means of awakening for the
eternal, then it does not harm at all; it helps.
That a miracle happens can be both for
help and for ruin; that the miracle does not happen
can be both for ruin and for help.
% --- p. 44 ---
\setcounter{footnote}{0}%
\opage{44}A human being is helped by God only when he is
helped unconditionally personally; but he is helped
unconditionally personally only when he is helped to
faith, in faith, by faith. Why has the
miracle happened? That I shall believe. Why
does the miracle not happen? That I shall believe. He
who denies the possibility of the miracle denies Providence;
he who assumes that God can help only
by a miracle also denies Providence.

As a sign of his divine omnipotence
the Son of Man performs the miracle:
he feeds the hungry. But --- is it not
wonderful? --- He who helps others by  % print: hiælper (misprint)
a miracle does not help himself by a miracle.
He who was able to feed five thousand
in the wilderness did not, however, turn stones into bread
to still his own hunger in the wilderness.

Here is seen the two-sidedness of Providence's help.
The solitary hungry one, who is tempted to a miracle,
and yet will not help himself by a miracle,
is nevertheless essentially helped: he has
the help personally in himself; he has the help
in the divine certainty that a human being
is not helped in one way, but in countless ways,
because ``man doth not live by bread alone,
% --- p. 45 ---
\setcounter{footnote}{0}%
\opage{45}but by every word that proceedeth out of the
mouth of God.''

The hungry multitude, which is helped
by the miracle, is nevertheless essentially helpless.
It hungers again the next day, and
seeks the Helper, not because it saw the sign,
but because it ate of the bread and was filled.
The help was offered, and yet it could not
help; because the hungry believed that bread alone
was at stake, because the multitude believed that
help was to be given temporally in the temporal.

With God all things are possible. In countless
ways Providence can help; and yet it is
all in vain when man will not
understand the eternal, personal help, and let
himself be helped. If help could be given temporally in the
temporal, then surely all the hungry would soon be filled,
all the naked clothed, and all the sorrows
of finitude ended; here are five thousand who
can bear witness that Providence both can and will
help; but here are also five thousand
who can bear witness that this is not the help.

With God all things are possible. In a supernatural
way Providence can reveal its
dominion over life and death. If the
% --- p. 46 ---
\setcounter{footnote}{0}%
\opage{46}healing of the sick and the resurrection of the dead
were the help itself: then omnipotence would surely quickly
heal all the sick, and open all graves; but
when it is a matter of the perishable personality's
transformation into imperishability, help can
be given only by inward healing, inward
awakening. That the miracle happens becomes, in
the revelation of Providence, precisely a proof that
the miracle in and for itself is not the help, and thereby
at the same time a proof that the means of help can just as well
consist in this, that no miracles happen.

With God all things are possible. As a sign
of his divine omnipotence the Son of Man
performs a wonder of awakening. But
--- is it not strange after all? --- He who
knows within himself that he can wake the dead man
from the grave, behold, he stands like a helpless one
grieved at the grave and weeps.

To be able to raise the dead man and yet stand
and weep at the grave looks like a contradiction. Yes, certainly, no one who will use his
understanding can deny that here there must be
a contradiction. But whose is the blame? who
is it that here causes the contradiction? Not
the Helper, but the many helpless, who by
% --- p. 47 ---
\setcounter{footnote}{0}%
\opage{47}him are to be helped. If the helpless
could at once see in what the help consists, then
there would be help enough here, then He who stands
at the grave could in a moment save all, and glorify
himself. But the help is misunderstood,
the means is transformed; what ought to have been for
the Helper's glorification brings precisely persecution
and suffering upon him, makes him, humanly
speaking, helpless; this is precisely the contradiction,
and a contradiction so sad that
the Son of God himself must weep over it.

``He hath saved others, himself he cannot
save.'' --- How there can be truth
in a blasphemy! For it is literally
true that he hath saved others, and it is
also literally true that he now cannot
save himself --- by the miracle.

Had the means of help been the help itself;
had lie and malice, violence and injustice
been capable of being annihilated by a miracle: then the
sacred history of Providence would not have become a history of
sufferings, but a history of glorification
for God and men. Now, on the contrary, since the miracle
is only a means, and the means is perverted,
help is given by the opposite means, so that
% --- p. 48 ---
\setcounter{footnote}{0}%
\opage{48}help is given in all inwardness by this, that
no miracle happens. In order to help personally
the personal Helper must himself become
like the helpless one. ---'' ``And will all then
be helped?''

With men this is impossible, but
with God all things are possible.

\afdeling[\rule{10em}{0.4pt}]{C.}{On Fate and Providence.}
\stykke{I.}{Fate Is Guilt.}

Because that which may be known of God is manifest to them, for God hath shewed it unto them, --- so that
they are without excuse.

He who would trust in the help of Providence
must do the will of Providence; he who
would do the will of Providence must know the thought of Providence. Who knows the thought of Providence?
Who does the will of Providence? The personal
truth, the true personality.

% --- p. 49 ---
\setcounter{footnote}{0}%
\opage{49}The relation between the One, who is the personal
truth, and the many, who through him
are to be made partakers of the truth, is not an
external historical relation between a solitary
individual who once existed, and
a multitude of other individuals who in time have
joined together about this One. Were it so,
there would be a difference of fate
between those who were so fortunate as to
enter into a personal relation to the extraordinary
individual, and the many, many who were
lawfully excused for not knowing him, because
they had never seen him, and never heard a
word about him.

The personal truth is the King of
truth; everyone who loves the truth hears
his voice, and everyone who hears his voice % print: Nøst (misprint)
is saved by Providence.

All fate is personal untruth; for fate
comes only by the personal itself
making itself into the untrue. All fate is essentially
guilt; for fate comes only by the personal wronging itself by
bowing under the impersonal. All fate is
essentially sin; for fate arises only
% --- p. 50 ---
\setcounter{footnote}{0}%
\opage{50}by man's denying God by
transforming him into something impersonal.

The immediate self-feeling of personality
is determined partly by the finite, partly
by the eternal. The self-feeling, immediately
determined by the finite, is egoism; determined
by the eternal, it is conscience.

The ideal of egoism is a mixture of the
personal with the impersonal; this mixture
is the veiled, in which the finite
personality hides its impotence from itself;
this veiled is the mythical. It is
not merely the idol that is mythical; the mythical
is in the idolater's impersonal personality.\footnote{The \textit{mythical} is not a determination of fantasy; it is an ethical determination. The mythical interpretation of the Old as well as of the New Testament is grounded in personal untruth, or, if one will, in impersonal truth.}

The ideal of conscience is a unity of revelation
and faith. In revelation the
egoistic impotence becomes manifest to itself in the face of
the personal omnipotence, and in faith man receives
the power of his true personality as a
power of God.

% --- p. 51 ---
\setcounter{footnote}{0}%
\opage{51}All ideals of egoism stand under fate;
all ideals of conscience are under
Providence.

Throughout the whole of history the struggle goes on between
personal truth and personal untruth,
between the ideals of egoism and of conscience. In the course of the struggle the race divides
asunder, and the movement parts into two currents:
the mythical, and that which in faith in revelation is \emph{religious}.

In the mythical direction the ideals of egoism
prevail amid ceaseless counteractions of conscience;
in the believing direction conscience prevails
amid ceaseless counteractions of
egoism. The separation is thoroughgoing, like a
separation between two opposed
spheres of spirit. Over both there is a fate,
for in both there is a guilt; but also over
both there is a Providence that has compassion, and
punishes the guilt.

For the people of the myths the punishment is fate;
for the people of revelation fate is punishment;
both are led to the goal: the personal truth
is the goal of Providence, and the goal is reached in
the fullness of time. 

% --- p. 52 ---
\setcounter{footnote}{0}%
\opage{52}In the fullness of time, then, the realized
ideal of conscience appears: \emph{saving}, for
the salvation of personality is the concern of conscience;
\emph{judging}, for every human being shall be judged
by his ideal. But the truth of the ideal is the King
of truth; every gleam of personal truth,
every spark of conscience that manifests itself
in the world is a spark of his flame, a
gleam of his light.

The personal truth is the King of
truth; everyone who loves the truth hears
his voice, and recognizes that he is the ideal itself.
But those who remain in untruth do not hear
his voice; they prefer to keep their own
self-made, distorted ideals, the one in nature,
the other in the written law. Now the
guilt becomes manifest. The offense now lies not
in the ideal's having been obscured and distorted  % print: deri,. (misprint)
in the course of the ages; for this distortion
is surely removed precisely by the King of truth
himself appearing in the world; no, the offense lies just
in this, that each of the parties in his own way wants to
keep the distorted ideal, the ideal that
contradicts conscience but flatters egoism.
Thus the misapprehension of
% --- p. 53 ---
\setcounter{footnote}{0}%
\opage{53}Providence becomes open rebellion against Providence;
such a rebellion is a blindness, the blindness
is fate, fate is guilt: here
no excuse holds.

If truth is posited as impersonal, then what the modern spirit asserts
is quite correct, namely that the history of the human race is sinless,
and that sins fall upon all individuals
without exception; but if the personal
truth is acknowledged, then the word of faith holds, that
the history of the race is a sinful history, and that
there is One in the race who is sinless.

The history of the world is in unceasing rebellion against
Providence; but the governance shows itself in this, that
the rebels are impotent. In order to see the rebellion
in its whole extent we must distinguish between
the noisy rebels and the silent rebels.

In a picture of the spirit the noisy
rebels are depicted from ancient days as the peoples that ``rage
and imagine a vain thing, as the mighty ones of the earth
who would break asunder the bands of
governance and strive against the King of Providence.'' Against
these the strong word is used, that He who
sitteth in the heavens shall laugh.

The punishment that is to strike these rebels
% --- p. 54 ---
\setcounter{footnote}{0}%
\opage{54}comes accordingly also in the shape of fate; for
the King of Providence shall'' ``break them with a rod of iron,
as a potter's vessel shall he dash them in
pieces.''

The noisy rebels are outwardly recognizable
precisely because they make noise; it is otherwise
with the silent ones. The silent rebels do not
boldly raise ideal against ideal; they do not lift
their heads so high that fate can historically
strike them: who, then, are the silent rebels?

Let us first confess that the silent rebellion smoulders
in us all, and further characterize the rebels
themselves according to the concept. They are surely all those who
trust in the assistance of Providence without ever
putting to themselves the serious question: what then, really, is the thought of Providence?
they are surely all those who confidently appeal to
truth and right, without however being willing to bow
unconditionally before the ideal of justice, and to hear
the King of truth; they are surely all those human beings,
high and low, rich and poor, gifted
and ungifted, who solely and alone, or at least chiefly
and most immediately, arrange their lives for
cozy well-being, agreeable security, finite
busyness, temporal activity for temporal
% --- p. 55 ---
\setcounter{footnote}{0}%
\opage{55}purposes, and for the rest are all virtue-adorned, in so far as
they are all adorned with the virtues most necessary for
earthly well-being, such as:
respectability, decency, industriousness,
sobriety, honesty in money matters, peaceableness,
obligingness, to which may further be
reckoned good-naturedness, affability, sociability,
etc., all these, who have
all these virtues and neither more nor less, all
these are neither more nor less than honnête,
decent, quiet, respectable rebels against the eternal
thought of Providence, against the King of truth
who humbled himself, and let the dispensations
come over him, that it might yet become manifest
what it is to know the thought of Providence
and to do the will of Providence.

\stykke{II.}{Temporal Conditions.}

The Lord maketh poor, and maketh rich:
he bringeth low, and lifteth up.

The inequality in human conditions is indeed
in a worldly sense to be regarded as a sorrowful
disproportion, an injustice of fate;
but in the sense of faith this inequality
% --- p. 56 ---
\setcounter{footnote}{0}%
\opage{56}is just a proof of the omnipresent justice of Providence.

In a temporal respect the rich man has an extraordinary  % print: Rige. (misprint)
favoring above the poor man, the
mighty above the lowly; in the
language of worldliness we must express this thus: Fate is he who maketh poor, and who
maketh rich; it is he who \emph{bringeth low}, and who lifteth up.

If, on the other hand, we see the matter in the light of faith, then
there is in this no injustice at all, no arbitrariness,
no partiality. That this is not
an arbitrary assertion but an eternal truth
follows directly from the opposition between the
worldly and the religious.

In relation to fate, that is, in worldliness,
man's significance is in every respect
made finite. A human being can be something great in the world; but all worldly greatness
is relative: even the greatest individuality
vanishes as an insignificance in the immeasurable
whole. Conversely a human being can be
something very lowly; but even the lowliest
is yet always Something. As now all greatness and all
lowliness in the worldly sense is piecework, so is
% --- p. 57 ---
\setcounter{footnote}{0}%
\opage{57}all of man's significance piecework: worldliness
is just the parceled-out human being.

It is not the same man who is at once
both great and lowly; rather the one is great,
the other lowly; it is not the same man who
is at once both rich and poor; rather the
one is rich, the other poor; indeed the same man
is not even always the same: formerly he was
mighty, now he is merely lowly; formerly he was
poor, now he is rich.

If now this worldly parceling-out of the human being
is held fast in its worldliness, then we
have fate. Every human being who has his
life in this parceling-out expresses personally the
impersonal truth: Fate is he who
maketh poor, and who maketh rich; it is he
who bringeth low, and who lifteth up.

In the face of God, on the other hand, man in
his finitude is ``dust and ashes,'' an indifferent
magnitude, a pure nothing; and yet man is
in God's thought something infinitely great, spirit of
God's spirit, God's image. Here the finite, the relative, falls away
on both sides,
and with it the parceling-out too falls away.

Here the same man is essentially always the
% --- p. 58 ---
\setcounter{footnote}{0}%
\opage{58}same. The same man who here is lowly is
also great, for his greatness issues from
his lowliness; the same man who here is poor
is also rich, for in his poverty there is born personally
a richness of eternity.

But now the piecework, the world and worldliness?
Worldliness must go; but the world
remains, for the world is good; the piecework remains,
for this piecework is very good --- in the eyes of
Providence.

The Lord maketh poor, and
maketh rich; he \emph{bringeth low}, and lifteth up.

All dispensations are personal tasks. By means of
the piecework the one task becomes
many tasks, and precisely in this difference
for each human being a really true and personal
task.

What in each dispensation is divine
help is at the same time temporal burden, and everything
that looks like temporal help is conversely
divine burden.

If we set the mighty man and the lowly man over against
each other, then the task comes in wholly opposite
dispensations, and, so to speak, under two 
% --- p. 59 ---
\setcounter{footnote}{0}%
\opage{59}disguises. Each of the two receives immediately
the personal impression of existence
that corresponds to his position: a human being is
something great; a human being is something lowly.

If now the solution of the task consisted merely
in an appropriation of the personal impression that a
human being is something great, and can accomplish something
great in the world: then the mighty man would be arbitrarily
favored, for that impression he receives
precisely from all sides at every moment that he lives.

If conversely the solution of the task consisted
in an immediate appropriation of the impression that
a human being is something lowly; then the lowly man
would be directly favored by his lowliness, for
the world will surely teach him how lowly he is.

But --- here the same man is essentially always
the same. The solution of the task therefore consists for both in a working-up and
inward appropriation of both impressions; thereby the conditions
become equal.

Can now the mighty man work up the impression
of worldly greatness in such a way that he --- not
merely for edification in a quiet hour, but in all
worldly surroundings, in real personality --- at
every moment transforms this impression into a religious
% --- p. 60 ---
\setcounter{footnote}{0}%
\opage{60}feeling of infinite lowliness before God; can the
mighty man, when thousands look up to him with admiration,
yet always see the eye that looks down upon him
testingly: then he solves a great task;
then he understands that temporal help is divine
burden; then he will also bear his burden,
remain in the height, and lift up his head as one
who is set in the height, elucidating the thought
that the Lord is he who lifteth up.

Next comes the lowly man, but --- how great,
how exalted is not his task! What he
is, what he feels, thinks and does, is after all
so insignificant, for the whole so indifferent,
that no one takes notice of it, because it is not
worth noticing at all. And yet this
insignificant man and all the insignificant things he does shall
be of infinite significance. For him who is not
even known in the world every moment shall be
just as important, just as exciting, just as decisive
as for him who is placed so high that
his beck moves a world. If now
the lowly man can solve this task, if he, the
unnoticed, can really appropriate the personal
impression that in every instant he is noticed by Providence: then he will truly understand that that which
% --- p. 61 ---
\setcounter{footnote}{0}%
\opage{61}is temporal burden is divine help;
then, if it is to be so, he will gladly remain
in lowliness as the oppressed: he is not, after all,
oppressed by fate; he expresses the thought
that the Lord is he who bringeth low.

If now we place the mighty man alongside the
lowly man, as they have both placed
themselves under Providence; then the first will surely not
despise the last, and the last not envy
the first; no, now they will understand one another
in all inwardness, and help one another with all their strength,
the mighty man by his might, the lowly man by
his lowliness.

``The Lord maketh poor, and maketh
rich.''

When the rich, in grateful acknowledgment of
the wisdom of Providence, impress this sentence
as a divine guarantee of the right of property,
while the poor, in impatience over the arbitrariness
of fate, rise up in order to come to the aid of wisdom:
then the waters roar, then
the rebels meet from both sides, the silent with
the noisy; but --- He that sitteth in the heavens
shall laugh.

It is in the common opinion an easy
% --- p. 62 ---
\setcounter{footnote}{0}%
\opage{62}matter to thank Providence for riches, but very
hard to have to thank it for poverty. This
is, however, not the divine opinion.
\emph{The truth} says: ``blessed are the poor'';
the truth says: ``it is easier for a camel
to go through the eye of a needle than for a
rich man to enter into the kingdom of God.'' If then
what \emph{the truth} says is true, then
riches must be a great and terrible burden;
but if riches are so great and terrible
a burden, then the rich man is surely not favored
by God above the poor man.

What shall the poor man do? Understand God's
dispensation, and transform the burden into help.

What shall the rich man do? Understand the dispensation,
and transform the burden into help. But
then the rich man is to do the very same as
the poor man. If now both are, each in his own
way, to do the same; then they also need mutual support, and
ought consequently to help one another; but above all the poor man ought to
help the rich man, for the poor man surely has
the less to bear.

% --- p. 63 ---
\setcounter{footnote}{0}%
\opage{63}The poor man who is helped by God in such a way
that he --- in personal truth really helps
the rich man, proves precisely thereby in deed that
the Lord is he who maketh poor.

``It is easier for a camel to go through
the eye of a needle than for a rich man to understand the thought
of Providence and to do the will of Providence''; riches
thus make the solution of the task so difficult that
the burden must almost be regarded as superhuman.
What, then, is the personal counsel of the truth
to the rich man? That the rich man as soon as possible lessen
his superhuman burden, and help
himself by helping the poor man. Only when the
rich man thus helps the poor man, only when
the rich man personally before God acknowledges that he,
every time he relieves the need of the
needy, essentially helps himself, does he express
the truth that the Lord is he who maketh
rich.

So the conditions are equal; so we are all,
though under many disguises, children of one
Father; so justice is divine,
as God is just.

% --- p. 64 ---
\setcounter{footnote}{0}%
\opage{64}\stykke{III.}{Spiritual Conditions.}

Every good and perfect gift cometh from above, from
the Father of lights, with whom is no variableness,
neither shadow of turning.

If anyone should ask outright whether the talents,
the geniuses, the heroes and the great ideas of the world
are not, in virtue of their higher nature, to be said in
quite a peculiar sense to be under Providence;
such a question cannot be answered
so long as it is held in a generality
that excludes the personal. Only when the
personal is taken in does the question
have meaning, and the answer will then depend
on the position that personality takes up
toward the ideas.

He who holds to the impersonal
truth, that is, in this case --- where the
personal is precisely the decisive thing --- to the
personal untruth, will then naturally maintain
that the ideas originate from the natural ground
of the world-reason, and that the higher gifts of spirit are consequently nothing other than pure gifts of nature.

He, on the other hand, who believes personally in the
personal truth will unconditionally hold fast the
% --- p. 65 ---
\setcounter{footnote}{0}%
\opage{65}presupposition that every good and perfect gift
cometh from above, from the Father of lights, that thus the ideas in history, like the birds of
the air, the geniuses in art, like
the lilies of the field, stand under God's fatherly
care.

The difference between the gifted and the ungifted
is in a certain way comparable to the difference
between the rich and the poor. The gifted
are indeed the spiritually rich, who impart to us of
their abundance what they can well spare.
The inequality, however, acquires a quite different significance in the spiritual than in the temporal.
Seen in the idea, nothing is so estimable
as a true aristocracy of spirit, and seen in the idea,
nothing is so contemptible as the mere aristocracy of money.

The spiritually gifted are destined to be
the benefactors, teachers, educators and
disciplinarians of the race; they are the servants of the idea, sent out
to serve those who are to be made partakers
of the riches of the ideas; but whichever of these
servants we may choose for an example, we
become convinced that every extraordinary
gift transforms itself into a divine burden,
% --- p. 66 ---
\setcounter{footnote}{0}%
\opage{66}when the gifted man is to express the thought
of Providence.

For God every human work, even the
greatest masterpiece, is nothing. For talent the
work is everything; for it is indeed ``the work that shall praise
the master.'' For God only the personal
ideal of conscience counts; in the relative spheres of spirit
the many impersonal ideals leading to the
execution of works count. If now the
gifted are to place themselves under Providence, then
each of them will have to live for
two ideals that counteract each other; for each
of them then receives, besides the ideal of conscience,
which makes the work a nothing, also
the ideal corresponding to his talent,
which makes the work everything. But to live for
two ideals that strive against each other is extraordinarily burdensome.

That the superiority of the gifted over the ungifted
is really a divine burden can be
shown in examples from the various spheres of spirit.

Suppose the brilliant artist places himself
under Providence beside the simple workman. They both then acknowledge the
% --- p. 67 ---
\setcounter{footnote}{0}%
\opage{67}ideal of conscience; they are thus equal in religious enthusiasm.
But the workman's work lies
so far below the ideal that it can in no way
become the object of tempting enthusiasm;
the artist's labor, on the other hand, is impossible without
a corresponding enthusiasm.

What a difference! While the workman only
worships one God, and has only one enthusiasm,
the artist almost gets two gods, in that he
gets two ideals, and is enthused in two heterogeneous
directions. While one will hardly find
any workman who thinks that his work
as work is to the glory of God (for it would
be strange if anyone asserted that he could
dig gravel to the glory of God, or slake lime to the
glory of God), the artist is incessantly tempted
to confound the demand of the ideal of art and that of the ideal of conscience.
Yes, if only talent could succeed
in transforming the worship of art
into a direct worship of God, then the
artist would be helped, then one could indeed
say that the gifted man was the favorite both of fortune and
of God. But it is impossible: aesthetic
pietism is rebellion against Providence.
No more can a picture be painted, be it
% --- p. 68 ---
\setcounter{footnote}{0}%
\opage{68}a Catholic picture of Mary or a Protestant
picture of Christ, to the glory of God, than an
overture can be composed to the glory of God, or
a comedy played to the glory of God.

All religious worship of art is refined idolatry.
In art the work is everything; in the eyes of
Providence the work is nothing. The separation of
sacred and secular art is so far from
marking a real difference between the religious
and the profane in the personal sense, that
the ``sacred'' art precisely brings with it the greatest
temptation to profane the sacred. The difference is, to be sure, drawn from the subject matter, but yet chiefly
from the inspiration of the talent. The profane
artist (the dramatist, the actor) is inspired
in a worldly way, and is not so easily tempted to
confuse his muse with an angel of Providence;
the ``sacred'' artist, on the other hand, gets his inspiration
through the sacred objects, his
enthusiasm for art is like an enthusiasm for God:
so it easily happens that the enthused man becomes
giddy, receives the prompting as a revelation
of God, falls worshipping upon his knees,
and becomes --- a refined idolater.

Every extraordinary talent is a good and
% --- p. 69 ---
\setcounter{footnote}{0}%
\opage{69}perfect gift; but it costs incredible exertion
to express the personal truth
that the gift cometh from above.

The not-gifted man, who eats his bread in the sweat of his
face, and sighs under the burdens of finitude,
has in these burdens a divine
help, and is not tempted by the relative ideals.
The gifted man, who is moved in mood and surrenders
himself in ideal productivity, must redouble his
exertion in order to preserve his personality.
The work, which is everything to him, must at the same time be to him
a nothing, an indifferent thing, a child's game, a vehicle
for the development of ethical powers. The
gifted man who can now do this: who can have
art as though he had it not, can love
its ideal as though he loved it not, because
he loves God above all things, can thank God
just as inwardly when the work utterly fails
as when it succeeds beyond measure, can hold
himself so unchanged to the ideal of conscience
that in his personality there is not so much as shadow
of turning, however fortune may change in
the world of art; yes, the gifted man who can
do this, he is truly a master over
many masters; he knows what it means
% --- p. 70 ---
\setcounter{footnote}{0}%
\opage{70}that the gift is more than a gift of fortune; he
expresses personally the truth that the gift
cometh from above, from the Father of lights, with
whom is no variableness, neither shadow of
turning.

A life consecrated to art always has
the possibility of much fortune and much misfortune, of
bitter pains and of rich enjoyments; this
holds directly, in so far as the gift is conceived as
a gift of nature. Thus abandoned to a fantastic
change of luck and ill luck, talent stands directly
under fate.

If, however, the gifted man will place himself under
Providence, then his life, regardless of
fortune and misfortune, is essentially consecrated to ethical exertion
and religious suffering. And what are
then all the quackeries of finitude about luck and
ill luck, recognition and non-recognition in
comparison with the inner personal exertion
of having to work for a goal, an
ideal, that is nevertheless a means, a nothing; of
having to be kindled by enthusiasm, and then again
to be enthused to quench the flame of enthusiasm!

If we now consider this, and bear in mind that the
children of art are the servants of the idea, sent out to
% --- p. 71 ---
\setcounter{footnote}{0}%
\opage{71}serve us; and consider what these servants
must sacrifice for our sake, bear and suffer for our
sake, then we surely have no cause to
envy them their spiritual superiority; we are
much more bound to serve them in turn in
everything we can, to meet them with gratitude,
and to honor them as the envoys of Providence,
provided they really (as we must believe)
strive in everything to express the thought
of Providence.

If we look about us in the kingdom of truth, we find
two wholly opposed spheres,
that of science and that of personality. Scientific
truth is in its essence impersonal;
personal truth is in its essence unscientific.

In science talent counts; it is again
the difference of finitude between the gifted and
the ungifted that proves that science belongs
to the world. In the sphere of personality this
difference is removed. The truth says: ``I thank
thee, O Father, Lord of heaven and earth,
that thou hast hid these things from the wise and prudent,
and hast revealed them unto babes'';
this says, after all, in other words: ``I thank
% --- p. 72 ---
\setcounter{footnote}{0}%
\opage{72}thee, Father, that thou hast made the personal
truth unscientific.''

The concepts fate and providence are themselves
unscientific concepts; their difference shows
itself only as personality steps in, and takes up
its position.

If personality bows under the scientific
truth, then the impersonal prevails,
and then comes fate. With this position
it is easy enough to dissolve the whole doctrine of faith
into untruth; but by this position personality
has at the same time essentially given itself up, and
finds itself in personal untruth. Scientific
truth has absolutely no guilt in
this untruth: the guilt lies in the personal
self.

If personality, on the other hand, asserts its original
right, and turns believing to the
object of faith: then it is a fairly easy matter to
get rid of the scientific objections; for
they fall away of themselves.

Let us then in the name of Providence acknowledge
our religious relation to God as a personal
relation; but let us then also, in the name of Providence,
resist the temptation to want
% --- p. 73 ---
\setcounter{footnote}{0}%
\opage{73}to have a ``\emph{science of faith},'' or in other
words: to want to have this personal element again interpreted
in the impersonal.

All that can be done scientifically in order
to clarify and secure the tasks of faith consists in
a complete and distinct separation of the boundaries of science
and of personality. Every concept of faith must therefore be presented
in such a way that it becomes recognizable precisely by its
sharp opposition to the concept of knowledge, and the
heterogeneity of the spheres demonstrated and elucidated
at all points.

Such categories as ``religious genius,''
``scientific gift of grace,'' belong to the offspring
of fancy, which are neither religious
nor scientific.

Genius is in its immediate form a gift
of nature, and no gift of nature can relate
itself directly as supporting the religious.

The religious moment is never scientific;
the scientific moment is never religious.

It goes with the man of science as with
the artist: if he would express the thought of Providence,
he is drawn in two opposite directions.
Science is his goal, his work, his
% --- p. 74 ---
\setcounter{footnote}{0}%
\opage{74}worldly ideal. Working for this goal,
he must use his whole power of consciousness in order to
remove everything personally disturbing and to forget himself
in the interest of the universal truth.
If he is then to return to the religious, he needs
conversely his whole power of consciousness in order to
collect himself personally, to rouse himself to an immediate
sensation of the punctual impression of actuality,
and to recognize that it is now not a matter of
this or that in general, but that it
concerns him, him himself personally before the personal God.

Scientificity does not further the religious,
it stupefies the religious. He who
sits in the deepest scientificity is as far
from the religious as he who lies
in the deepest sleep.

Here too it appears that worldly advantage
is divine burden. It is far easier amid the occupations, sorrows and
joys of everyday life to work inwardly in the religious than
to have the religious in a study. You
can carry a millstone, and have the religious;
but no thinker on earth can have
the religious at the moment he unties a
% --- p. 75 ---
\setcounter{footnote}{0}%
\opage{75}scientific knot. A woman can spin at her
spinning wheel, and have the religious; but a man,
were he ever so believing, cannot reason
scientifically about God's existence, and at the same
time have the religious.

That all science in and for itself, like all
that is worldly in and for itself, leads away from God,
lies in this, that all science leads away from the
personal. When the personal, the punctual in the ideal of conscience,
is first duly weakened,
the ideal of science can then also easily transform
itself into the Godhead itself. Hence such
phrases are heard as: ``Science brings us
nearer to God; the cultivation of science is in a higher
sense the worship of God.''

But is science, then, not a good and perfect
gift? Yes, truly! but it costs
incredible exertion, if the existence of the knower
is to express the personal truth that the
gift cometh from above, from the Father of lights,
with whom is no variableness, neither shadow
of turning.

The superiority of spirit is man's suffering:
He who hath the eternal Spirit has suffered for
us all.

% --- p. 76 ---
\setcounter{footnote}{0}%
\opage{76}If we survey the distribution of the spiritual
gifts as a whole, it shows itself, from two different
sides: worldly unjust, divinely
just.

The spirit of worldliness, the spirit of the relative ideals,
of finite works and strivings,
is impersonal, and for that very reason unjust. The gifts of nature
are distributed indiscriminately like the gifts of
fortune: the one is gifted, the other limited;
the one completes a great work, and
becomes something great; the other accomplishes nothing,
and remains in insignificance. So the labors are divided:
the gifted man gets the great labor,
and is at an advantage; the ungifted man must stand back.
So the laborers contend over the labor;
so the laborers are worn out and consumed as
means, and the works, the final goals, remain
behind as the yield of history, as Babylon's
rubble, as Egypt's pyramids, as memorials
of the masses of personal powers
that in the course of the ages have been sacrificed to fate.

The spirit of Providence, of the ideal of conscience, of the eternal
purpose is personal, and for that very reason
just. In this kingdom of spirit each
human being receives his gift as personal task.
% --- p. 77 ---
\setcounter{footnote}{0}%
\opage{77}He who received the greatest power received also the
greatest burden. Here the gift corresponds exactly to
the burden, the work to the responsibility, the reward to the exertion;
here the greatest receives no more,
and can receive no more, than the least; for
the least as well as the greatest receives,
each according to his exertion in faith, the whole
personal truth.

\stykke{IV.}{The One Thing Needful.}

But without faith it is impossible to please God; for he
that cometh to God must believe that
he is, and that he is a rewarder of them that
seek him.

``Nature must help itself'': this rule,
which surely holds much validity in medicine,
seems almost capable of being transferred to ethics,
when one sees everything worldly. For is it
not strange how the personal natural dispositions
can mock upbringing and cultivation;
how a good nature can bring itself up
in spite of bad examples, and a bad nature harden
itself against the best examples; how an independent
nature, now pliant, now contentious, can
% --- p. 78 ---
\setcounter{footnote}{0}%
\opage{78}maintain its individuality under strong external
pressure, while the weak, dependent
nature cannot by any supporting
surroundings whatever be helped over its innate
immaturity.

In the worldly sense there is here of course, in
this as in every other temporal inequality, an
arbitrariness of chance, an injustice of fate;
but in the divine sense
fate is guilt, and the guilt is man's
own.

That this is so can indeed be
heard, when those concerned begin to complain
of fate.

If the weak man complains that he must stand back
for the strong man; he may well be right
in the temporal; for in everything temporal the
strong nature accomplishes far more than the weak. But
before God the weak man accuses himself; for before
God faith counts as the one thing needful, and
in faith the difference is removed.

In order to believe, the strong man must become weak,
like the weakest; for God is strong only in the
weak. If there were here any finite difference,
then the advantage would surely rather be on
% --- p. 79 ---
\setcounter{footnote}{0}%
\opage{79}the side of the weaker; but here there is no difference:
natural weakness and natural
strength vanish alike in the weakness that
alone counts before God. So equality before
God is posited in the inward, and all the inequality that belongs
to the outward, the whole outward-going activity,
is then to be regarded as the vanishing,
which does not concern the personal self.

If it is a restless, passionate mind, tormented in inner
distraction, that would compare
itself with the deliberate, calm individuality, harmoniously
disposed by nature, and complain
of its fate: then this complaint too will be
heard by God as a self-accusation; for faith
is demanded, and in faith the difference is removed.

The naturally calm man must be shaken and disquieted
to the utmost, in order to reach faith; and
the naturally restless man finds rest and peace
when he comes to faith. Faith is not
a little rest and a little unrest, or sheer peace in some
and sheer unpeace in others; but faith is everywhere
and in all the eternal's unrest in all natural rest,
the eternal's peace in all natural unpeace. The
inequality in the mixture of peace and unpeace, that
as a consequence of the individualities' 
% --- p. 80 ---
\setcounter{footnote}{0}%
\opage{80}natural diversity remains, belongs to the temporal
world, the temporal struggle, and is subordinate to the personal
self.

If the gentler, softer nature would complain
in comparison with the harder and hardier, that the latter
finds it far easier to bear
the dispensations of life; then this complainer too is heard
by God as one who accuses himself.
Existence may be hard and severe; the nature
that is hardy has in the temporal the
unmistakable advantage that it hardens itself more easily
against the blows of fate than the unfortunate man
who received the softer, gentler mind. But in faith
the difference is removed. The task of faith is not too
severe for anyone, for it is infinitely mild; it
is not too mild for anyone, for it is infinitely severe.
The heavenly Father is not like an earthly
father a little more compliant toward the compliant, a little
more hard toward the hard; no, the God
of Providence is essentially equally gracious and equally severe
toward all natures.

As there can be an obstinacy in
hardness, so there can also be an obstinacy in whimpering softness: the holy
severity consumes all natural obstinacy;
% --- p. 81 ---
\setcounter{footnote}{0}%
\opage{81}but, when the divine fire has consumed
the human obstinate element, it transforms itself
into mildness, into infinite mercy. Therefore
no human being shall know his faith by this, that
he has severity alone, and no one by this,
that he has mildness alone; but amid all the
differences of individualities everyone can
know his faith by this, that it always has severity
in mildness and mildness in severity.

At times we human beings complain that faith
is too unearthly, too spiritual, too overstrained,
so that it does not really permit us to look at the
temporal, to rejoice in the purely human
and to work in the present. It is natural
that this complaint arises in everyone
who to any degree exerts himself with the task
of faith; let us then complain, the one about this,
the other about that, if only we know that
before God we accuse ourselves.

One thing is needful; this One is faith, in
faith unconditional acknowledgment of the demands of
conscience, in faith unconditional humbling
under the ideal of personality, in faith
unconditional obedience to the laws of the ideal, in faith
unconditional trust in the promises of the ideal; when
% --- p. 82 ---
\setcounter{footnote}{0}%
\opage{82}a human being in inward earnestness strives to hold fast
this One: then let him live in the world as
he can and will, solitary or in company, fasting
in the desert or enjoying in society: there shall
be blessing in it all, and it shall
all turn out to his good.

But without faith it is impossible.

The temporal as temporal is not to a certain
degree indifferent to God; it is infinitely indifferent.
Temporal want, temporal pain, temporal
joy, temporal enjoyment, all this is in itself
nothing; here no pity of finitude counts,
no whimpering and no lamenting, all
temporal things go and shall go according to the order of nature:
man in flesh and blood is ``dust
and ashes,'' a nothing for God, as the world in all
its glory and all its misery is a
nothing for God, the Almighty. But, precisely because
the temporal in its temporality is infinitely indifferent,
it is, as means and condition for
man's personal education, of infinite
validity. Providence has not, like Odin, only
one eye; Providence has millions of eyes, an
eye with every human being, an eye in every conscience.
In the throng of the generations, in the press of
% --- p. 83 ---
\setcounter{footnote}{0}%
\opage{83}events no soul is forgotten, and
no accident is accidental.

The powers of the world make a din, the Almighty is silent;
the powers of the world arm themselves and defend their
right, the Almighty is so still and calm!
Let those who can be content with a little right and
a little wrong, a little guilt and a little fate, interpret
this as they will: every human being who in the name of conscience
would be saved from the violence of fate
must acknowledge that Providence is revealed,
that the knot is untied, that there is help for everyone
who will let himself be helped.

But without faith it is impossible; for faith
in Providence is the help of Providence.

\begin{center}---\end{center}

\end{document}
